% U3 WS2 -- Type A/B: gas laws, KMT, real gases. Blocks 3-5.
\documentclass{apchem-worksheet}

\apcunit{3}
\apcunittitle{Properties of Substances and Mixtures}
\apcdoctitle{Worksheet 2 \textbullet\ Gases}
\apcsubtitle{CED 3.4--3.6 \textbullet\ kelvin every time}

\begin{document}
\wsheader[Convert every temperature to kelvin before substituting. Show the
rearranged equation before plugging in numbers.]

\problem[3.4]{Ideal gas calculations. $R = \SI{0.08206}{\liter\atm\per\mole\per\kelvin}$.}
\begin{parts}
  \item Volume of \SI{3.00}{\mole} at \SI{2.00}{\atm} and
        \SI{127}{\celsius}: \workspace[2cm]
        \keyonly{\begin{apcanswer}$T = \SI{400.}{\kelvin}$;
        $V = (3.00)(0.08206)(400.)/2.00 = \SI{49.2}{\liter}$\end{apcanswer}}
  \item Moles in \SI{5.00}{\liter} at \SI{1.20}{\atm} and
        \SI{27}{\celsius}: \answerblank[3cm]{\SI{0.244}{\mole}}
  \item Pressure of \SI{0.100}{\mole} in \SI{2.00}{\liter} at
        \SI{300.}{\kelvin}: \answerblank[3cm]{\SI{1.23}{\atm}}
\end{parts}

\problem[3.4]{A gas occupies \SI{6.0}{\liter} at \SI{300}{\kelvin} and
\SI{1.0}{\atm}.}
\begin{parts}
  \item Volume at \SI{600}{\kelvin}, constant pressure:
        \answerblank[2.6cm]{\SI{12}{\liter}}
  \item Volume at \SI{3.0}{\atm}, constant temperature:
        \answerblank[2.6cm]{\SI{2.0}{\liter}}
  \item A student converts \SI{300}{\kelvin} to \SI{27}{\celsius} and uses
        27 in the ratio. Explain the error's effect.
        \answerlines[2]{Ratios require absolute temperature; using Celsius
        would predict a 22-fold volume change instead of doubling, because
        Celsius has an arbitrary zero.}
\end{parts}

\problem[3.4]{A \SI{10.0}{\liter} vessel holds \SI{0.30}{\mole} \ce{N2},
\SI{0.50}{\mole} \ce{O2}, and \SI{0.20}{\mole} \ce{Ar} at
\SI{300.}{\kelvin}.}
\begin{parts}
  \item Total pressure: \answerblank[3cm]{\SI{2.46}{\atm}}
  \item Mole fraction of \ce{O2}: \answerblank[3cm]{0.50}
  \item Partial pressure of \ce{O2}: \answerblank[3cm]{\SI{1.23}{\atm}}
  \item If all the \ce{O2} were removed, what happens to the partial
        pressure of \ce{N2}? \answerlines[2]{Unchanged --- each gas exerts
        pressure independently of the others.}
\end{parts}

\problem[3.5]{State which KMT assumption each observation depends on:}
\begin{parts}
  \item Gas pressure rises when the container shrinks.
        \answerblank[7.3cm]{constant random motion / collisions}
  \item All gases at the same $T$ have equal average kinetic energy.
        \answerblank[5.5cm]{$KE \propto T$}
  \item Gas particles don't slow down over time in a sealed flask.
        \answerblank[5.5cm]{elastic collisions}
\end{parts}

\problem[3.5]{\ce{He} and \ce{CO2} are held at the same temperature.}
\begin{parts}
  \item Compare their average kinetic energies.
        \answerblank[6.4cm]{equal --- KE depends only on $T$}
  \item Compare their average speeds, with reasoning.
        \answerlines[2]{He moves much faster: $u_{\text{rms}} =
        \sqrt{3RT/M}$, and He's molar mass is 11 times smaller.}
  \item Sketch both Maxwell--Boltzmann curves on one set of axes and label
        them. \workspace[3.4cm]
        \keyonly{\begin{apcanswer}He: broad, flat, peak far right.
        \ce{CO2}: narrow, tall, peak left. Equal areas.\end{apcanswer}}
\end{parts}

\problem[3.6]{Under which conditions does a real gas behave most ideally,
and why? Name both failed assumptions in your answer.
\answerlines[3]{High temperature and low pressure. At low pressure
particles are far apart, so their own volume is negligible compared with the
container; at high temperature they move fast enough that intermolecular
attractions barely affect collisions.}}

\problem[3.6]{For \ce{CH4} at \SI{200}{\kelvin}, a plot of $PV/nRT$ vs $P$
first dips below 1, then rises above 1.}
\begin{parts}
  \item Cause of the dip: \answerblank[6.6cm]{attractions reduce collision
        force}
  \item Cause of the rise: \answerblank[7.2cm]{particle volume becomes
        significant}
  \item Would \ce{He} show a larger or smaller dip? Why?
        \answerlines[2]{Smaller --- He's dispersion forces are extremely
        weak, so attractions barely reduce the measured pressure.}
\end{parts}

\begin{spiralstrip}{Unit 3 \textbullet\ blocks 1--2}
  \item Strongest IMF in \ce{CH3CH2OH}: \answerblank[3.5cm]{hydrogen bonding}
  \item Solid type of diamond: \answerblank[3.5cm]{covalent network}
  \item Higher vapor pressure: water or acetone (bp \SI{56}{\celsius})?
        \answerblank[2.5cm]{acetone}
  \item Melting a molecular solid breaks:
        \answerblank[3.5cm]{IMFs, not bonds}
  \item Higher bp: \ce{Ar} or \ce{Ne}? \answerblank[2cm]{\ce{Ar}}
\end{spiralstrip}

\end{document}
